\section{Overview [F]}

In this part, explain own implementation for the Script Assembler in detail. Aside from \textit{Syntax Highlighting}, everything was written in the \textit{Python} programming language.

The following chapters belong to this part:
\begin{itemize}
    \item Chapter~\ref{chap:preprocessor}: \textit{Preprocessor}. This chapter goes over the \textit{Preprocessor}, which converts significant indentation of a \textit{modifier}'s \textit{block body} into a reserved keyword that is understood by the \textit{Parser}. In Chapter~\ref{chap:ebnf-problem}, the need for this preprocessor is explained in detail.
    \item Chapter~\ref{chap:parser}: \textit{Parser}. This chapter goes over the \textit{Parser}, which uses the \textit{Lark} library in combination with a modified form of our EBNF from Chapter~\ref{lst:our_ebnf} to recursively parse an input script to an output script and brackets. 
    \item Chapter~\ref{chap:option-modifier-format}: \textit{Option/Modifier Format}. Here, our implementation for handling the respective keywords within the parser is explained. 
    \item Chapter~\ref{chap:trp-integration}: \textit{TRP-Integration}. Since our project is used within the \gls{trp} project, we explain the changes we made to the existing project to integrate our Script Assembler.
    \item Chapter~\ref{chap:syntax-highlighting}: \textit{Syntax Highlighting}. Since the existing \gls{trp} project's \textit{Angular} frontend uses a code editor known as \textit{Monaco Editor}, we implemented \textit{Syntax Highlighting}, in other words text colorization based on keyword syntax, which we go over in this chapter.
    \item Chapter~\ref{chap:extensibility}: \textit{Extensibility}. In this chapter, we explain how the Script Assembler may be extended with additional keywords, give examples to keywords that may be useful in the future, and discuss the limits of our parser. 
    \item Chapter~\ref{chap:database}: \textit{Database}. Here, we explain the \textit{Mockup Database}, which we built to allow the project to run independently to the proprietary \textit{TestRail} system and to aid with testing.
\end{itemize}

[where to find py]